% GENERATED FILE. Edit canonical lessons, then run npm run generate:books.
% canonical-source-hash: fnv1a64:c27458f3c978907c
% canonical-lessons: TE-C458-sarduko, TE-C458-alavatupadu, TE-C458-domatera, TE-C458-marugudoddi, TE-C458-taguniru

\chapter{Settling In}
\label{ch:te-settling-in}
\hlchaptermodality{458}

\begin{chapteropening}
\textbf{By the end of this chapter:} \emph{I can say to settle in, to get used to, a mosquito net, a toilet and drinking water.}

Complete the last lesson of chapter 458: I can say to settle in, to get used to, a mosquito net, a toilet and drinking water.
\end{chapteropening}

\section[\texorpdfstring{\te{సర్దుకో} (sardukō)}{సర్దుకో (sardukō)}]{\te{సర్దుకో} (sardukō) — to settle in; to make do}

\label{lesson:TE-C458-sarduko}



\begin{quote}
Before the new one: say the Telugu for an elevator, then the Telugu for to stay overnight.
\end{quote}

\subsection*{You'll want to know: \te{సర్దుకో}}
\textbf{\te{సర్దుకో}} — \emph{sardukō} — \textquotedblleft{}to settle in; to make do\textquotedblright{}.

\begin{quote}
\te{కొత్త} \te{గదిలో} \te{రెండు} \te{రోజుల్లో} \te{సర్దుకున్నాను}.
\end{quote}

\begin{quote}
\emph{kotta gadilō reṇḍu rōjullō sardukunnānu.}
\end{quote}

\begin{quote}
I settled into the new room in two days.
\end{quote}

Past (I …): \te{సర్దుకున్నాను}, \emph{sardukunnānu}.

\te{సర్దు}, to arrange, + -\te{కో}: to arrange things for oneself, to make do with what there is.

\begin{cousinweb}
\textbf{\te{సర్దుకో}} (\emph{sardukō}) is built on \te{సర్దు}, \textquotedblleft{}to arrange, to tidy\textquotedblright{}, whose origin is not settled, with the ending -\te{కో}, \textquotedblleft{}for oneself\textquotedblright{}, from \te{కొను}, \textquotedblleft{}to take\textquotedblright{} (DEDR 2151).
\end{cousinweb}

\begin{grammarlens}[title={What you've built}]
One more word for getting around.
\end{grammarlens}

\subsection*{Your turn}
\begin{itemize}
\raggedright
  \item \emph{Say it:} \emph{sardukō}
  \item \emph{Say it:} \emph{sardukō}, once more
  \item \emph{Say it:} \emph{kotta gadilō reṇḍu rōjullō sardukunnānu.} — I settled into the new room in two days.
  \item \emph{Recall:} say the Telugu for an elevator, then the Telugu for to stay overnight, then say \emph{sardukō} again
\end{itemize}

\begin{tcolorbox}[breakable,colback=teal!4,colframe=teal!35!black,title={Before you move on}]
What does \te{సర్దుకో} mean? (\textquotedblleft{}To settle in; to make do\textquotedblright{}.) Say it once more.
\end{tcolorbox}
\section[\texorpdfstring{\te{అలవాటుపడు} (alavāṭupaḍu)}{అలవాటుపడు (alavāṭupaḍu)}]{\te{అలవాటుపడు} (alavāṭupaḍu) — to get used to}

\label{lesson:TE-C458-alavatupadu}



\begin{quote}
Before the new one: say the Telugu for to stay overnight, then the Telugu for to settle in.
\end{quote}

\subsection*{You'll want to know: \te{అలవాటుపడు}}
\textbf{\te{అలవాటుపడు}} — \emph{alavāṭupaḍu} — \textquotedblleft{}to get used to\textquotedblright{}.

\begin{quote}
\te{చల్లని} \te{వాతావరణానికి} \te{నెమ్మదిగా} \te{అలవాటుపడ్డాను}.
\end{quote}

\begin{quote}
\emph{callani vātāvaraṇāniki nemmadigā alavāṭupaḍḍānu.}
\end{quote}

\begin{quote}
I slowly got used to the cold weather.
\end{quote}

Past (I …): \te{అలవాటుపడ్డాను}, \emph{alavāṭupaḍḍānu}.

What you get used to takes -\te{కి}: \te{వాతావరణానికి}. \te{అలవాటు}, a habit, + \te{పడు}.

\begin{cousinweb}
\textbf{\te{అలవాటుపడు}} (\emph{alavāṭupaḍu}) is native Dravidian: \te{అలవాటు}, \textquotedblleft{}a habit, a practice\textquotedblright{} (DEDR 249), with \te{పడు}, \textquotedblleft{}to undergo\textquotedblright{} (DEDR 3853).
\end{cousinweb}

\begin{grammarlens}[title={What you've built}]
One more word for getting around.
\end{grammarlens}

\subsection*{Your turn}
\begin{itemize}
\raggedright
  \item \emph{Say it:} \emph{alavāṭupaḍu}
  \item \emph{Say it:} \emph{alavāṭupaḍu}, once more
  \item \emph{Say it:} \emph{callani vātāvaraṇāniki nemmadigā alavāṭupaḍḍānu.} — I slowly got used to the cold weather.
  \item \emph{Recall:} say the Telugu for to stay overnight, then the Telugu for to settle in, then say \emph{alavāṭupaḍu} again
\end{itemize}

\begin{tcolorbox}[breakable,colback=teal!4,colframe=teal!35!black,title={Before you move on}]
What does \te{అలవాటుపడు} mean? (\textquotedblleft{}To get used to\textquotedblright{}.) Say it once more.
\end{tcolorbox}
\section[\texorpdfstring{\te{దోమతెర} (dōmatera)}{దోమతెర (dōmatera)}]{\te{దోమతెర} (dōmatera) — a mosquito net}

\label{lesson:TE-C458-domatera}



\begin{quote}
Before the new one: say the Telugu for to settle in, then the Telugu for to get used to.
\end{quote}

\subsection*{You'll want to know: \te{దోమతెర}}
\textbf{\te{దోమతెర}} — \emph{dōmatera} — \textquotedblleft{}a mosquito net\textquotedblright{}.

\begin{quote}
\te{దోమతెర} \te{లేకుండా} \te{రాత్రి} \te{నిద్ర} \te{పట్టలేదు}.
\end{quote}

\begin{quote}
\emph{dōmatera lēkuṇḍā rātri nidra paṭṭalēdu.}
\end{quote}

\begin{quote}
Without a mosquito net I could not sleep at night.
\end{quote}

\te{దోమ}, a mosquito, + \te{తెర}, a screen, a curtain.

\begin{cousinweb}
\textbf{\te{దోమతెర}} (\emph{dōmatera}) is native Dravidian: \te{దోమ}, \textquotedblleft{}a mosquito\textquotedblright{} (DEDR 2991), and \te{తెర}, \textquotedblleft{}a screen, a curtain\textquotedblright{} (DEDR 3244).
\end{cousinweb}

\begin{grammarlens}[title={What you've built}]
One more word for getting around.
\end{grammarlens}

\subsection*{Your turn}
\begin{itemize}
\raggedright
  \item \emph{Say it:} \emph{dōmatera}
  \item \emph{Say it:} \emph{dōmatera}, once more
  \item \emph{Say it:} \emph{dōmatera lēkuṇḍā rātri nidra paṭṭalēdu.} — Without a mosquito net I could not sleep at night.
  \item \emph{Recall:} say the Telugu for to settle in, then the Telugu for to get used to, then say \emph{dōmatera} again
\end{itemize}

\begin{tcolorbox}[breakable,colback=teal!4,colframe=teal!35!black,title={Before you move on}]
What does \te{దోమతెర} mean? (\textquotedblleft{}A mosquito net\textquotedblright{}.) Say it once more.
\end{tcolorbox}
\section[\texorpdfstring{\te{మరుగుదొడ్డి} (marugudoḍḍi)}{మరుగుదొడ్డి (marugudoḍḍi)}]{\te{మరుగుదొడ్డి} (marugudoḍḍi) — a toilet}

\label{lesson:TE-C458-marugudoddi}



\begin{quote}
Before the new one: say the Telugu for to get used to, then the Telugu for a mosquito net.
\end{quote}

\subsection*{You'll want to know: \te{మరుగుదొడ్డి}}
\textbf{\te{మరుగుదొడ్డి}} — \emph{marugudoḍḍi} — \textquotedblleft{}a toilet\textquotedblright{}.

\begin{quote}
\te{బస్టాండులో} \te{శుభ్రమైన} \te{మరుగుదొడ్లు} \te{ఉన్నాయి}.
\end{quote}

\begin{quote}
\emph{basṭāṇḍulō śubhramaina marugudoḍlu unnāyi.}
\end{quote}

\begin{quote}
There are clean toilets at the bus stand.
\end{quote}

Plural \te{మరుగుదొడ్లు}. A polite, public word; at home people also say \te{బాత్రూం}.

\begin{cousinweb}
\textbf{\te{మరుగుదొడ్డి}} (\emph{marugudoḍḍi}) is native Dravidian: \te{మరుగు}, \textquotedblleft{}concealment, cover, a screen\textquotedblright{}, which DEDR 4760 lists as Telugu \emph{maṟũgu}, and \te{దొడ్డి}, \textquotedblleft{}a backyard, an enclosure\textquotedblright{} (DEDR 3485): a screened-off place.
\end{cousinweb}

\begin{grammarlens}[title={What you've built}]
One more word for getting around.
\end{grammarlens}

\subsection*{Your turn}
\begin{itemize}
\raggedright
  \item \emph{Say it:} \emph{marugudoḍḍi}
  \item \emph{Say it:} \emph{marugudoḍḍi}, once more
  \item \emph{Say it:} \emph{basṭāṇḍulō śubhramaina marugudoḍlu unnāyi.} — There are clean toilets at the bus stand.
  \item \emph{Recall:} say the Telugu for to get used to, then the Telugu for a mosquito net, then say \emph{marugudoḍḍi} again
\end{itemize}

\begin{tcolorbox}[breakable,colback=teal!4,colframe=teal!35!black,title={Before you move on}]
What does \te{మరుగుదొడ్డి} mean? (\textquotedblleft{}A toilet\textquotedblright{}.) Say it once more.
\end{tcolorbox}
\section[\texorpdfstring{\te{తాగునీరు} (tāgunīru)}{తాగునీరు (tāgunīru)}]{\te{తాగునీరు} (tāgunīru) — drinking water}

\label{lesson:TE-C458-taguniru}



\begin{quote}
Before the new one: say the Telugu for a mosquito net, then the Telugu for a toilet.
\end{quote}

\subsection*{You'll want to know: \te{తాగునీరు}}
\textbf{\te{తాగునీరు}} — \emph{tāgunīru} — \textquotedblleft{}drinking water\textquotedblright{}.

\begin{quote}
\te{గ్రామంలో} \te{తాగునీరు} \te{దొరకడం} \te{కష్టంగా} \te{ఉంది}.
\end{quote}

\begin{quote}
\emph{grāmaṁlō tāgunīru dorakaḍaṁ kaṣṭaṅgā undi.}
\end{quote}

\begin{quote}
It is hard to get drinking water in the village.
\end{quote}

\te{తాగు}, to drink, + \te{నీరు}, water: water for drinking. \te{నీళ్ళు} is the everyday plural.

\begin{cousinweb}
\textbf{\te{తాగునీరు}} (\emph{tāgunīru}) is native Dravidian: \te{తాగు}, \textquotedblleft{}to drink\textquotedblright{} (DEDR 3174), and \te{నీరు}, \textquotedblleft{}water\textquotedblright{} (DEDR 3690).
\end{cousinweb}

\begin{grammarlens}[title={What you've built}]
One more word for getting around.
\end{grammarlens}

\subsection*{Your turn}
\begin{itemize}
\raggedright
  \item \emph{Say it:} \emph{tāgunīru}
  \item \emph{Say it:} \emph{tāgunīru}, once more
  \item \emph{Say it:} \emph{grāmaṁlō tāgunīru dorakaḍaṁ kaṣṭaṅgā undi.} — It is hard to get drinking water in the village.
  \item \emph{Recall:} say the Telugu for a mosquito net, then the Telugu for a toilet, then say \emph{tāgunīru} again
\end{itemize}

\begin{tcolorbox}[breakable,colback=teal!4,colframe=teal!35!black,title={Before you move on}]
What does \te{తాగునీరు} mean? (\textquotedblleft{}Drinking water\textquotedblright{}.) Say it once more.
\end{tcolorbox}
